\documentclass[11pt]{article}
\usepackage[margin=0.9in]{geometry}
\usepackage{amsmath,amssymb,amsthm,booktabs}
\usepackage[T1]{fontenc}
\usepackage{lmodern,microtype}
\usepackage[colorlinks=true,linkcolor=blue,urlcolor=blue]{hyperref}
\newtheorem{theorem}{Theorem}[section]
\newtheorem{lemma}[theorem]{Lemma}
\newtheorem{proposition}[theorem]{Proposition}
\newtheorem{corollary}[theorem]{Corollary}
\theoremstyle{remark}
\newtheorem{remark}[theorem]{Remark}
\newcommand{\R}{\mathbb R}
\newcommand{\C}{\mathbb C}
\newcommand{\norm}[1]{\left\|#1\right\|}
\newcommand{\Herm}{\operatorname{Herm}}
\title{Toward a four-dimensional Yang--Mills construction\\
\large Notebook 06: direct flat Maxwell bounds and a compact lift}
\author{Research working document}
\date{16 October 2026}
\setlength{\emergencystretch}{5em}
\begin{document}
\maketitle

\begin{abstract}
This note proves bounds for the flat quadratic effective action that
are uniform in the number of averaging levels. Curvature averaging
is a contraction in the correctly scaled volume norms. A finite-block
Poincar\'e argument therefore supplies an axial gap. An explicit compact
right inverse to straight-string averaging supplies the complementary
energy upper bound. A gauge adjustment shows that the rooted and
straight constraints have the same energy infimum. Together with the
actual averaged-path gauge bridge of Notebook 05, these results give
both bounds for the flat reduced Gaussian precision, subject to the
stated identification with the constrained-minimizer Hessian. Complex
background responses and covariance locality remain separate problems.
\end{abstract}

\noindent\textbf{Status.}
The proofs have passed separate internal checks of the averaging,
finite-block gap, compact lift, periodic counting, and variational
transfer. No external certification or novelty claim is made.
The statements below are about an explicitly
defined flat quadratic model. Section~\ref{sec:match} gives the precise
additional identification needed for the deposited Yang--Mills
Gaussian; no nonlinear renormalization theorem or continuum limit is
inferred from these bounds.

\section{Normalized averaging commutes with curvature}

Let fine and coarse spacings be $h$ and $H=Mh$, with $M$ odd.
Use centered $M$-blocks $B_M(y)$, one positive bond orientation,
and one plaquette orientation $\mu<\nu$. Potentials are physical
coordinates: the group link would be $\exp(ihA_\mu(x))$.
Define
\[
 (d_hA)_{\mu\nu}(x)=h^{-1}
 [A_\nu(x+he_\mu)-A_\nu(x)
  -A_\mu(x+he_\nu)+A_\mu(x)],
\]
\[
 (Q_M^{\rm st}A)_\mu(y)=M^{-d-1}
 \sum_{x\in B_M(y)}\sum_{r=0}^{M-1}A_\mu(x+rhe_\mu),
\]
\[
 (Q_M^{(2)}F)_{\mu\nu}(y)=M^{-d-2}
 \sum_{x\in B_M(y)}\sum_{r,s=0}^{M-1}
 F_{\mu\nu}(x+rhe_\mu+she_\nu).
\]
The volume norm is
$\norm F_h^2=h^d\sum_x\sum_{\mu<\nu}|F_{\mu\nu}(x)|^2$,
and analogously for one-forms.

\begin{proposition}\label{prop:cochain}
On the infinite lattice or a periodic lattice tiled by complete blocks,
\begin{equation}\label{eq:cochain}
 d_HQ_M^{\rm st}=Q_M^{(2)}d_h,\qquad
 \norm{Q_M^{(2)}F}_H\le\norm F_h.
\end{equation}
The one-form average is also a contraction in the volume norms.
\end{proposition}

\begin{proof}
Expand a coarse plaquette curl and translate the block base points.
The four strings are the boundary of an $M$ by $M$ fine plaquette
square. Discrete Stokes contributes $h\sum d_hA$, and the coarse
curl contributes $H^{-1}=(Mh)^{-1}$. Together with $M^{-d-1}$,
this gives exactly $M^{-d-2}$.

Each row of $Q_M^{(2)}$ is a positive average of $M^{d+2}$ terms.
For a fixed fine plaquette and each pair $(r,s)$, the base point lies
in exactly one block. Its total column weight is therefore
$M^2M^{-d-2}=M^{-d}$. Jensen's inequality, followed by
$H^dM^{-d}=h^d$, proves the contraction. Periodic preimages are
counted with multiplicity, so small tori cause no exception.
Use $M$ string indices in place of $M^2$ surface indices for the
one-form contraction.
\end{proof}

The physical derivative $h^{-1}$ is essential in this computation.
Unscaled plaquette sums would change the power of $M$.
Straight-string averages satisfy the exact semigroup identity
\[
 Q_N^{\rm st}Q_M^{\rm st}=Q_{NM}^{\rm st}.
\]
Indeed a centered $NM$-block decomposes uniquely into centered
$M$-blocks with intermediate base coordinates, and a string index
decomposes uniquely as $r_M+Mr_N$. The two bijections are independent.
The same proof gives the two-form semigroup identity.

\subsection{The rooted contour average}

Let $T_hA(y,x)$ be the physical rooted path integral, including the
factor $h$, and let $a$ average it over the block. The linearization
of the rooted group average is
\begin{equation}\label{eq:rooted}
 Q_M^\flat=Q_M^{\rm st}-D_HaT_h.
\end{equation}
The initial and final rooted contours give the second term; the
reference coarse string cancels with the final reference transporter.
This is the spacing-restored identity of Notebook 05. Since
$d_HD_H=0$, the curvature identity \eqref{eq:cochain} also holds
with $Q_M^\flat$. Iterating gives
\begin{equation}\label{eq:iterated-contraction}
 \norm{d_1Q_kA}_1\le\norm{d_\eta A}_\eta,
 \qquad \eta=L^{-k},
\end{equation}
for the composed rooted derivative. No equality of its one-form
operator with $Q_{L^k}^{\rm st}$ is used.

The corrected operator in Raw I, lines 24355--24375, has a different
flat derivative: all four correction brackets vanish at the identity,
leaving $Q^{\rm st}$ exactly. Its straight semigroup is valid. The
rooted operator of CMP109 differs by the gradient in
\eqref{eq:rooted}; the two should not be identified merely by notation.

\section{A compact right inverse with bounded energy}

Set coarse spacing to one and fine spacing to $\eta=1/M$. We now
construct a right inverse $R_M$ to $Q_M^{\rm st}$ with constants
independent of $M$. This section uses infinite or periodic complete
block lattices. Restriction to an arbitrary boundary is discussed
below.

For $M=2h_0+1\ge3$, put
\[
 m_1=\frac{M^2-1}{4M^2},\qquad
 m_2=\frac{M^2-1}{12M^2},\qquad
 v=m_2-m_1^2
   =\frac{(M^2-1)(M^2+3)}{48M^4}.
\]
Define a longitudinal sequence, zero outside $0\le s<M$, by
\begin{equation}\label{eq:psi}
 u_s=(s-h_0)/M,\qquad
 \psi_M(s)=\alpha+\beta|u_s|,
 \qquad \alpha=m_2/v,\quad\beta=-m_1/v.
\end{equation}
For the normalized mean $E_M=M^{-1}\sum_{s=0}^{M-1}$,
\begin{equation}\label{eq:psi-moments}
 E_M\psi_M=1,\qquad E_M(|u|\psi_M)=0,
 \qquad E_M\psi_M^2=\alpha\le4.
\end{equation}
These follow by substituting $E_M|u|=m_1$, $E_Mu^2=m_2$.
In particular the signed sequence is even about $h_0$.

For transverse coordinates define the hat sequence
\begin{equation}\label{eq:phi}
 \phi_M(r)=\frac{M}{h_0+1}
             \left(1-\frac{|r|}{h_0+1}\right)
 \quad (|r|\le h_0),
 \qquad\phi_M(r)=0\quad\hbox{otherwise}.
\end{equation}
Direct summation gives
\begin{equation}\label{eq:phi-bounds}
 M^{-1}\sum_r\phi_M(r)=1,\quad
 0\le\phi_M\le2,\quad M^{-1}\sum_r\phi_M(r)^2\le2,
\end{equation}
\begin{equation}\label{eq:phi-gradient}
 M^{-1}\sum_r|M[\phi_M(r+1)-\phi_M(r)]|^2
 =\frac{16M^3}{(M+1)^3}\le16.
\end{equation}
There are $M+1$ nonzero differences, all with absolute value
$M/(h_0+1)^2$ before multiplying by $M$.

For a coarse bond based at $j\in\mathbb Z^d$ in direction $\mu$,
define the fine bond basis, writing $z=Mx\in\mathbb Z^d$, by
\[
 \Psi_{\mu,j}(z)=\psi_M(z_\mu-Mj_\mu)
                \prod_{\nu\ne\mu}\phi_M(z_\nu-Mj_\nu),
\]
and set
\begin{equation}\label{eq:lift}
 (R_MB)_\mu(z/M)=\sum_jB_\mu(j)\Psi_{\mu,j}(z).
\end{equation}
The sums are local. Each basis reads only its own coarse bond and
is supported in the union of the two endpoint blocks. Periodize the
basis for a torus. For $M=1$, use the identity lift.

\begin{theorem}[Compact bounded-energy lift]\label{thm:lift}
For $M$ odd and $d\ge2$,
\begin{equation}\label{eq:lift-bounds}
 Q_M^{\rm st}R_M=I,\qquad
 \norm{R_MB}_\eta^2\le2^{d+1}\norm B_1^2,
\end{equation}
\begin{equation}\label{eq:lift-energy}
 \norm{d_\eta R_MB}_\eta^2
 \le C_d\norm B_1^2,
 \qquad C_d=2^{d+6}(d-1).
\end{equation}
The bounds hold in infinite volume and on periodic lattices, including
small tori with periodic multiplicities.
\end{theorem}

\begin{proof}
The averaging separates into coordinates. Transversely, the supports
of $\phi_M(r-Mj)$ are disjoint blocks and their block averages are
Kronecker deltas by \eqref{eq:phi-bounds}.
Longitudinally, the straight-string kernel in a row based at $Mj$
is the triangle
\[
 M^{-1}\left(1-\frac{|s-Mj-h_0|}{M}\right)_+,
\]
on its integer support. On the support of the basis with index $j$,
the target row thus averages $(1-|u|)\psi_M$, which is one by
\eqref{eq:psi-moments}. The adjacent rows average the positive and
negative halves of $|u|\psi_M$. Both are zero by evenness and the
second moment condition. Every other row has disjoint support.
This proves $QR=I$; periodization preserves the identity.

For each orientation the tensor supports tile the fine sites without
overlap. The squared norm of a basis is bounded by
$4\,2^{d-1}=2^{d+1}$, using the volume-weighted tensor product
of \eqref{eq:psi-moments} and \eqref{eq:phi-bounds}.
This proves the first norm bound.

Only transverse derivatives $D_{\eta,\nu}A_\mu$, $\nu\ne\mu$,
enter the curl. Thus no derivative of the signed longitudinal
sequence is needed. For one basis, its transverse derivative has
squared norm at most $4\,16\,2^{d-2}=2^{d+4}$.
At a given fine bond, at most two differentiated translates overlap.
Cauchy--Schwarz therefore gives
\[
 \norm{D_{\eta,\nu}(R_MB)_\mu}_\eta^2
 \le2^{d+5}\norm{B_\mu}_1^2.
\]
Finally sum
$|D_\mu A_\nu-D_\nu A_\mu|^2\le
2|D_\mu A_\nu|^2+2|D_\nu A_\mu|^2$
over $\mu<\nu$. Each component occurs $d-1$ times, proving
\eqref{eq:lift-energy}. The same overlap counting on the periodic
cover proves the bound for small tori. The identity lift at $M=1$
obeys the displayed conservative constants directly.
\end{proof}

For example, at $M=3$, $\psi=(0,3,0)$ and
$\phi=(3/4,3/2,3/4)$. At larger $M$, the longitudinal basis has
negative portions; that cancellation is what makes it an exact
right inverse. Positivity of the lift is neither asserted nor needed.

For supported coarse fields on the full lattice this lift remains
supported in the corresponding double blocks. A coarse bond incident
to a bounded region can have one endpoint block outside it. The
unmodified basis then needs a collar. The following separate
construction removes that need for the specified incident-bond
boundary convention.

\subsection{An exact incident-bond boundary lift}

Let $S$ be any collection of coarse sites and let the fine region
be the union of their complete centered $M$-blocks. Retain bonds
with at least one endpoint in the corresponding region; set every
other coarse or fine coordinate to zero. The constraint is the
full straight-string average of this zero-extended field, with no
truncation or renormalization of boundary rows. The energy is the
full curl norm of the zero extension.

For $M=2h_0+1\ge3$, put
\[
 a_M=\frac{2M^2}{(h_0+1)(h_0+2)}\le8.
\]
Replace the longitudinal profile on an outgoing coarse bond (initial
site in $S$, final site outside) by
\[
 \psi_M^{\rm out}(s)=a_M[1-3(h_0-s)/M]
       \quad(0\le s\le h_0),
\]
zero elsewhere. For an incoming positive bond use its reflection
\[
 \psi_M^{\rm in}(s)=a_M[1-3(s-h_0)/M]
       \quad(h_0\le s<M),
\]
zero elsewhere. Keep the original profile on bonds with both endpoints
inside. The interface bond $s=h_0$ is retained in either boundary
case, and is included only once in the selected profile.

For $p(r)=a_M(1-3r/M)$, $0\le r\le h_0$, direct summation gives
\[
 \frac1M\sum_rp(r)=1,\qquad
 \frac1M\sum_r\frac rM p(r)=0,\qquad
 \frac1M\sum_rp(r)^2=a_M\le8.
\]
Indeed $\sum r=h_0(h_0+1)/2$ and
$\sum r^2=h_0(h_0+1)M/6$ prove the first two equalities;
multiplying $p=a_M(1-3r/M)$ by $p$ proves the third.
The target triangular averaging row therefore remains one. One
neighboring row vanishes by the first-moment identity, and the other
has disjoint support. Every further row vanishes. Thus every selected
profile still maps to its exact single coarse bond.

An outgoing profile uses only fine bonds inside the initial block
and its crossing interface bond. An incoming profile uses the
crossing bond and bonds inside the final block. With the unchanged
transverse hats, every used fine bond has an endpoint in the allowed
region, even when $S$ has holes. Tensoring and summing as in
\eqref{eq:lift} defines $R_M^S$ with
\begin{equation}\label{eq:boundary-lift}
 Q_M^{\rm st}R_M^S B=B,\qquad
 \norm{R_M^S B}_\eta^2\le2^{d+2}\norm B_1^2,
\end{equation}
\[
 \norm{d_\eta R_M^SB}_\eta^2
 \le2^{d+7}(d-1)\norm B_1^2.
\]
The constants are independent of the region and refinement. The
proof of Theorem~\ref{thm:lift} applies with longitudinal squared
norm eight in place of four. In particular, differentiate each
fixed basis before summing: even when adjacent source bonds use
different profile types, only two transverse derivative supports
overlap. No derivative of a position-dependent profile choice is
introduced. The same argument works on periodic covers, and $M=1$
uses the identity. In four dimensions the boundary energy constant
is $6144$.

The identity holds on all coarse rows, including the excluded ones.
Indeed an excluded row's strings lie in its two exterior endpoint
blocks and read no allowed fine bond. Thus full-row and incident-row
constraints agree for this convention. The one-step deposited flat
convention matches: Raw II, lines 9032--9039, uses incident bonds;
the positivity setup beginning at line 9743 uses complete blocks;
lines 9779--9791 use $Q_1$, whose flat derivative is $Q^{\rm st}$.

This proves the boundary upper bound for that exact linear problem.
It does not identify an arbitrary multilevel retained-history
minimizer, establish a boundary covariance estimate, or cover a
convention freezing crossing bonds. Nonlinear transfer still needs
the corresponding constraint and gauge match of
Section~\ref{sec:match}.

\subsection{An independent periodic Fourier upper bound}

There is a separate check of scale uniformity. For a coarse Fourier
mode $k\in[-\pi,\pi]^d$, embed the same physical frequency in the
fine torus. Let
\[
 r_M(k)=\frac{\sin(k/2)}{M\sin(k/(2M))},\qquad r_M(0)=1.
\]
For centered blocks the averaging multiplier is
\[
 q_\mu(k)=e^{i(M-1)k_\mu/(2M)}r_M(k_\mu)
                              \prod_{\nu=1}^d r_M(k_\nu).
\]
Since $|r_M(k)|\ge2/\pi$ on this zone, divide the component
coefficient by $q_\mu(k)$ to obtain a right inverse. Fine-volume
Parseval equals coarse Parseval for the embedded frequencies. The
exterior derivative multiplier has squared norm at most
$\sum_\nu|M(e^{ik_\nu/M}-1)|^2\le d\pi^2$. Hence
\[
 \inf_{Q_M^{\rm st}A=B}\norm{d_\eta A}_\eta^2
 \le d\pi^2(\pi/2)^{2d+2}\norm B_1^2.
\]
For real $B$, taking the real part of a complex lift preserves the
constraint and cannot increase energy, covering Nyquist conventions.
Unlike \eqref{eq:lift}, this Fourier lift is not compactly supported.

\section{The flat variational form and rooted constraints}

Define, first as a standalone quadratic model,
\begin{equation}\label{eq:qmin}
 q_M(B)=\inf_{Q_M^{\rm st}A=B}\norm{d_\eta A}_\eta^2,
 \qquad\eta=M^{-1}.
\end{equation}
The constraint is nonempty by Theorem~\ref{thm:lift}. On a finite
torus this is a nonnegative quadratic form: it is least squares
on a nonempty affine space. Equations \eqref{eq:cochain} and
\eqref{eq:lift-energy} give directly
\begin{equation}\label{eq:energy-sandwich}
 \norm{d_1B}_1^2\le q_M(B)\le C_d\norm B_1^2.
\end{equation}
It is invariant under adding a coarse gradient. Indeed
$Q_M^{\rm st}D_\eta\lambda=D_1\overline\lambda$, where
$\overline\lambda$ is the scalar block average; this average is
onto, and gradients have zero curl energy.

The composed rooted constraint has the same infimum. In one step,
\[
 Q_M^\flat D_\eta\lambda=D_1(\lambda\hbox{ at centers}),
 \qquad
 Q_M^{\rm st}D_\eta\lambda=D_1\overline\lambda.
\]
Equation \eqref{eq:rooted} and induction therefore give
\[
 Q_k^{\rm rooted}A=Q_{L^k}^{\rm st}A+D_1\Xi_k(A)
\]
for a linear scalar map $\Xi_k$. If the straight constraint equals
$B$, extend the coarse-center values $-\Xi_k(A)$ to a fine site
field and add its gradient. Nested center restriction corrects the
rooted constraint to $B$, without changing curl energy. Conversely,
the scalar block average has a right inverse (blockwise constants),
so a gradient also corrects a rooted-feasible field to a
straight-feasible one. Both infima are equal. This uses no bound
on $\Xi_k$ and proves no bound on the gauge-adjusted potential or
its cutoff coordinates.

\section{A finite-block axial Poincar\'e inequality}

Fix the one-step blocking factor $L\ge3$, independent of $M=L^k$.
On the unit lattice let $\Omega$ consist of the $3^d$ adjacent
$L$-blocks indexed by $\{-1,0,1\}^d$. Include every bond and
plaquette internal to this rectangular grid. Let $T_{0,\Omega}$
contain the root-to-site integrals for a chosen coordinate-ordered
tree in every block, and let $Q_{{\rm st},\Omega}$ contain the
straight-string rows between all adjacent block centers inside
$\Omega$. Each such string stays in its two endpoint blocks.
Form the completely specified finite matrix
\[
 M_L=d_\Omega^*d_\Omega
       +T_{0,\Omega}^*T_{0,\Omega}
       +Q_{{\rm st},\Omega}^*Q_{{\rm st},\Omega},
 \qquad \sigma(d,L)=\lambda_{\min}(M_L).
\]

\begin{lemma}\label{lem:poincare}
One has $\sigma(d,L)>0$. Every global field with
$T_0B=0$ and $Q_1B=Q_L^{\rm st}B=0$ satisfies
\begin{equation}\label{eq:axial-gap}
 \norm{d_1B}_1^2\ge\gamma_{\rm loc}\norm B_1^2,
 \qquad \gamma_{\rm loc}=\sigma(d,L)/3^d>0.
\end{equation}
The constant is independent of volume, averaging depth, and color
dimension.
\end{lemma}

\begin{proof}
If the matrix quadratic form vanishes, all local plaquette curls
vanish. Path independence on a rectangular grid gives $B=D\varphi$;
interchanging adjacent coordinate steps is justified by each zero
plaquette curl. The tree equations force $\varphi$ to be constant
within every block. A straight-string row equals the difference of
the two block constants divided by $L$. The block graph is connected,
so all constants agree, and $B=0$. Thus $M_L$ has trivial kernel,
which proves strict positivity of its smallest eigenvalue.

Restrict a globally constrained field to each cube $\Omega_y$
translated by an aligned block center $y\in L\mathbb Z^d$.
Its local tree and averaging rows vanish, giving
\[
 \sum_{b\subset\Omega_y}|B_b|^2
 \le\sigma(d,L)^{-1}\sum_{p\subset\Omega_y}|(d_1B)_p|^2.
\]
Every global bond occurs in at least one cube, and each plaquette
in at most $3^d$ cubes. Summing proves \eqref{eq:axial-gap}.
On a torus pull the field back to the ordinary cubes in the covering
lattice. Even when a cube wraps, the same proof and counting with
multiplicity apply. In infinite volume, use square-summable fields
and monotone convergence.
\end{proof}

No numerical computation of $\sigma$ is needed for strict positivity.
The value may be small as a function of fixed $d,L$; its independence
of the ultraviolet scale is the relevant conclusion. Together with
\eqref{eq:energy-sandwich}, the lemma proves the full-field axial
bound $q_M(B)\ge\gamma_{\rm loc}\norm B^2$ directly.

\section{Both bounds for the flat exponential-gauge precision}

Use the averaged path map $T$, rooted constraint $Q^\flat$, and
pivot insertion $K_\flat$ from Notebook 05. Set
\[
 d_0=2\sqrt d,\quad
 t_0=\sqrt{d(L-1)(L^d-1)},\quad
 r=1+2d(L-1)/L,
\]
\[
 K_*^2=1+L^{2d}(2d+2)r^2,
 \qquad
 \gamma_*=
 \left((1+d_0t_0)^2/\gamma_{\rm loc}+d_0^2\right)^{-1}.
\]
The constraint-compatible gauge bridge gives, on $\ker Q^\flat$,
\[
 q_M(B)+\norm{TB}^2\ge\gamma_*\norm B^2.
\]
In fact the lower estimate needs only $q_M(B)\ge\norm{d_1B}^2$:
the maps $B\mapsto B_T\mapsto B_0$ of Notebook 05 change fields
by gradients, so their curls agree. Apply Lemma~\ref{lem:poincare}
to $B_0$, then $\norm{B_T}\le(1+d_0t_0)\norm{B_0}$, and finally
weighted Cauchy--Schwarz to $B=B_T+D(TB)$. Thus this particular
lower bound does not need a separate gauge-invariance hypothesis
on the effective quadratic form.
The direct upper bound \eqref{eq:energy-sandwich} applies on all
bond fields, so also
\[
 q_M(B)+\norm{TB}^2\le(C_d+t_0^2)\norm B^2.
\]
With $B=K_\flat z$, $\norm{K_\flat}\le K_*$ and $z$ a
coordinate restriction of $B$, the reduced real precision satisfies
\begin{equation}\label{eq:precision-sandwich}
 \gamma_*I\le A_{M,0}\le B_*I,
 \qquad B_*=(C_d+t_0^2)K_*^2.
\end{equation}
All constants depend only on fixed block geometry and dimension
in the displayed normalization. In particular,
\[
 B_*^{-1}I\le C_{M,0}=A_{M,0}^{-1}\le\gamma_*^{-1}I.
\]
If only the Wilson quadratic action is multiplied by $c>0$ while
the gauge penalty stays $\norm{TB}^2$, the appropriate constants are
\[
 \gamma_*(c)=
 \left((1+d_0t_0)^2/(c\gamma_{\rm loc})+d_0^2\right)^{-1},
 \qquad B_*(c)=(cC_d+t_0^2)K_*^2.
\]
Rescaling both terms together instead rescales the whole precision.

For later use, if the actual complex symmetric precision is
$A=A_{M,0}+V$ with $\norm V\le\gamma_*/2$, then
\[
 \Herm A\ge(\gamma_*/2)I,\qquad
 \norm A\le B_*+\gamma_*/2,
\]
\[
 \Herm A^{-1}\ge
 \frac{\gamma_*/2}{(B_*+\gamma_*/2)^2}I,
 \qquad\norm{A^{-1}}\le2/\gamma_*.
\]
This last implication is elementary; the actual uniform perturbation
estimate is not proved here. The spectral bounds alone do not give
exponential covariance decay. Notebook 07 supplies that additional
flat spatial estimate by an exact Fourier quotient and complex-pole
cancellation. Localized background dependence remains a separate input.

\section{Identification with the Yang--Mills flat Hessian}\label{sec:match}

For the Wilson action in physical normalization,
\[
 S_\eta(U)=\eta^{d-4}\sum_p[1-\Re\operatorname{tr}U(\partial p)],
\]
and $U_b(t)=\exp(i\eta tA_b)$, stationarity at the identity and
the first plaquette phase $i\eta^2t(d_\eta A)_p$ give
\[
 D^2S_\eta(1)[A,A]=\eta^d\sum_p|(d_\eta A)_p|^2.
\]
The trace fixes the Lie-algebra norm. A different fixed normalization
of the action multiplies its quadratic energy bounds by a positive
factor; the preceding paragraph explains the gauge-penalty adjustment.
Colors decouple at quadratic order.

Let $U_k(V)$ be the local full-block constrained minimizer with
$U_k(1)=1$ in its chosen gauge, and assume the branch is $C^2$
on a neighborhood of that datum. Write
\[
 U_k(e^{itB})=\exp(i\eta[tL_kB+O(t^2)]).
\]
The constraint gives $Q_kL_kB=B$. Since the fine action has zero
gradient, its effective Hessian is
\[
 q_{{\rm eff},k}(B)=\norm{d_\eta L_kB}_\eta^2
                  \ge\norm{d_1B}_1^2.
\]
The lower bound uses the derivative, constraint, and action expansion;
the $C^2$ hypothesis makes it a statement about the classical effective
Hessian. Mere differentiability would instead give the expansion
$tL_kB+o(t)$ and the corresponding quadratic Peano coefficient.
To identify the Hessian with the variational form
$q_M$ and obtain the upper bound, also use minimization against
every tangent-feasible curve. At each fixed finite regulator, the
linear average is onto by the lift and gradient adjustment above.
For the actual analytic constraint, the finite-dimensional implicit
function theorem therefore produces a feasible curve for every
tangent-feasible field, provided the constraint map is defined on
an open neighborhood of the identity. Comparing its action with
the local minimizer and dividing by $t^2/2$ gives
\[
 q_{{\rm eff},k}(B)=
       \inf_{Q_kA=B}\norm{d_\eta A}_\eta^2=q_{L^k}(B).
\]
The reverse inequality uses the feasible derivative $L_kB$ itself.
If the chosen chart imposes an additional gauge, one must supply
its constraint-preserving gauge equivalence before using every
such competitor. The equality is not obtained from stationarity
alone. No uniform nonlinear implicit-function radius is claimed
by this finite-regulator argument.

CMP109, (0.21)--(0.22) and (2.6)--(2.11), identifies the effective
action and the quadratic Gaussian in these terms. The deposited
Raw II small-field step around lines 73870--74020 records the same
minimizer variables and elimination. The first derivative of the
group average must be checked from its contour definition, as done
above; gauge covariance alone would not determine it. The corrected
and rooted conventions have equal flat energy infima only through
the explicit gradient adjustment proved here.

Notebook 15 now supplies the analytic ordinary-constraint local
minimum and its comparison modulo gauge on the homogeneous small
global potential class. Notebook 17 checks the preceding tangent
comparison and deposited coordinate match explicitly. Thus
\eqref{eq:precision-sandwich} applies to that homogeneous flat
source Gaussian. This does not apply automatically to domains
with holes or mixed large-field
histories, or identify a freely prescribed current with a physical
action derivative. It does not transport the sharp cutoff under
an untracked coordinate change.

For the quadratic cutoff coefficient evaluated at $(U,J)=(1,0)$,
this flat reference is the relevant starting point. What remains
is a uniform weighted bound on its actual complex perturbations
and source responses. Notebook 07 proves the homogeneous flat spatial
decay, exact compatible-periodic volume identification, and a free
refinement limit with $O(M^{-2})$ weighted error. These results remove
the flat spectral and spatial estimates from the outstanding inputs;
they do not complete the background neighborhood or a four-dimensional
Yang--Mills theory.

\begin{thebibliography}{9}
\bibitem{repo} \href{https://github.com/ethanwsussman/Yang--Mills-LLM-experiment/tree/e392ef45013a581664bab76b90867170b30fe10e}
{Pinned Yang--Mills-LLM-experiment account}, commit
\texttt{e392ef4501\allowbreak 3a581664ba\allowbreak b76b908671\allowbreak 70b30fe10e}.
\bibitem{balaban109} T. Ba\l aban,
\emph{Renormalization group approach to lattice gauge field theories I},
Commun. Math. Phys. 109 (1987), 249--301.
\href{https://scispace.com/pdf/renormalization-group-approach-to-lattice-gauge-field-3ygai0vxt1.pdf}
{Copy of the original article}.
\end{thebibliography}
\end{document}
